%% =====================================================================
%% Methodological Pitfalls in Machine-Learning Solar Forecasting
%% IEEE conference format (IEEEtran). Compiles on Overleaf as-is:
%%   New Project -> Upload Project -> select this .tex -> compile with pdfLaTeX.
%%
%% BEFORE SUBMISSION:
%%   - Fill in the affiliation/department/city placeholders in \author.
%%   - Re-run the experiments and confirm every number in Section V.
%%   - Verify page numbers and DOIs in the bibliography against the
%%     publisher's record.
%% =====================================================================

\documentclass[conference]{IEEEtran}
\IEEEoverridecommandlockouts

\usepackage{cite}
\usepackage{amsmath,amssymb}
\usepackage{booktabs}
\usepackage{graphicx}
\usepackage{url}
\usepackage{balance}

\def\BibTeX{{\rm B\kern-.05em{\sc i\kern-.025em b}\kern-.08em
    T\kern-.1667em\lower.7ex\hbox{E}\kern-.125emX}}

\begin{document}

\title{Quantifying Three Evaluation Pitfalls in\\
Machine-Learning Solar Forecasting:\\
Split Optimism, Interval-Label Phase Error,\\
and Prediction-Interval Miscalibration}

\author{
\IEEEauthorblockN{
T.~Harshini~Sowmya,
M.~Lavanya~Durga,
P.~J.~S.~Siddhartha,\\
M.~Ajay~Prakash,
P.~Maruthi~Sreeram,
and~K.~Surya~Prakash}
\IEEEauthorblockA{
\textit{Department of Computer Science and Engineering}\\
\textit{\textless Institution Name\textgreater}\\
\textless City\textgreater, India\\
Team CSE-D\_13}
}

\maketitle

% =====================================================================
\begin{abstract}
Machine-learning solar forecasting is routinely reported with
coefficients of determination above~0.9, yet several widely used
evaluation practices inflate those figures in ways that are invisible
in the published result. We reconcile six recent solar-forecasting
studies into a single reproducible platform and use it to quantify
three such effects on identical data, models and random seeds. First,
random train/test splitting of an autocorrelated hourly series
understates RMSE by 11.6\,\% (74.68 to 66.03~W/m\textsuperscript{2}) and
overstates $R^{2}$ by 0.027 relative to a chronological split---the
comparison that one of the reconciled studies explicitly names as
future work. Second, we determine empirically that hourly reanalysis
irradiance is labelled at the \emph{end} of its averaging interval;
evaluating solar position at the raw label rather than the interval
midpoint introduces a systematic phase error yielding 264 spurious
clear-sky exceedances per year, and correcting it raises our
data-quality score from 97.7\,\% to 99.9\,\%. Third, uncalibrated
quantile gradient boosting attains only 66.3\,\% empirical coverage for
a nominal 80\,\% prediction interval; conformalized quantile regression
restores coverage to 80.2\,\% and to within two points in every weather
regime. We additionally show that permutation feature importance is
contingent on target choice, which reconciles two apparently
contradictory published importance rankings, and that a scenario
response can conceal two large opposing pathways that nearly cancel.
None of these findings requires metered plant output: each is a
controlled comparison in which only the disputed methodological choice
varies. We argue that reporting them should be routine, and release the
platform and its reproducibility manifest to make doing so
inexpensive.
\end{abstract}

\begin{IEEEkeywords}
solar irradiance forecasting, photovoltaic systems, ensemble learning,
time-series cross-validation, conformal prediction, reproducibility,
ERA5 reanalysis
\end{IEEEkeywords}

% =====================================================================
\section{Introduction}

Solar generation is variable on every timescale that matters to an
operator, and the machine-learning literature addressing that
variability has grown quickly. Recent studies report strong
goodness-of-fit: $R^{2}$ near or above~0.9 is common, and ensemble
methods---random forests, boosted trees, and stacked
combinations---are consistently among the best performers
\cite{P2,P4,P5}.

The difficulty is that a headline $R^{2}$ is a function not only of the
model but of a set of evaluation choices that are frequently made in
passing and reported in a single clause. Whether night hours are
retained, whether the split respects chronology, whether the target was
statistically truncated, and whether a stated prediction interval
actually covers what it claims are all decisions that move the reported
number substantially. Because they are rarely varied within a single
study, their magnitude is difficult for a reader to judge.

This paper takes six recent solar-forecasting studies
\cite{P1,P2,P3,P4,P5,P6}, implements their common methodological core
in one platform, and then uses that platform to hold everything
constant except one disputed choice at a time. The result is a set of
effect sizes rather than a set of recommendations.

\subsection{Contributions}

\begin{enumerate}
  \item \textbf{Split optimism, measured.} With identical model, seed
    and data, random hour-level splitting understates RMSE by
    11.6\,\% and overstates $R^{2}$ by $+0.027$ against a chronological
    split. Blocked random splitting by whole days recovers about half
    the gap (Section~\ref{sec:f2}).

  \item \textbf{The interval-label convention, determined
    empirically.} Reanalysis irradiance is an interval average while
    solar position is instantaneous. We sweep candidate offsets over a
    full year and show the representative instant lies 30~minutes
    \emph{before} the label, eliminating clear-sky exceedances entirely
    (Section~\ref{sec:f1}). We are not aware of this convention being
    established by measurement elsewhere in the applied
    solar-ML literature, and a defect it exposed in our own code is
    reported alongside it.

  \item \textbf{Interval miscalibration, measured and corrected.}
    A nominal 80\,\% interval from plain quantile gradient boosting
    covered 66.3\,\% of observations. Conformalized quantile regression
    \cite{romano2019} restores 80.2\,\% coverage overall and holds
    within two points across clear, cloudy and precipitation regimes
    (Section~\ref{sec:f3}).

  \item \textbf{Target-dependent feature importance.} Grouped
    permutation importance ranks solar geometry third when the target
    is the clear-sky index but dominant when the target is raw power,
    which reconciles \cite{P5} with \cite{P2} rather than adjudicating
    between them (Section~\ref{sec:f4}).

  \item \textbf{A reproducible platform.} All comparisons run from one
    codebase against a keyless public data source, with a
    reproducibility manifest, so a reader can repeat them on their own
    coordinates.
\end{enumerate}

A deliberate non-contribution: we make \emph{no} claim of
state-of-the-art forecast accuracy. Section~\ref{sec:limits} states why
our setup cannot support one, and why the contributions above do not
require it.

% =====================================================================
\section{Related Work}

\subsection{Targets and models}

The six reconciled studies do not share a prediction target. Four
predict irradiance---global horizontal irradiance (GHI) in
W/m\textsuperscript{2} \cite{P2,P3,P4}, or global solar radiation
\cite{P4}---while the remainder predict PV power directly
\cite{P1,P5,P6}. The two are related by a deterministic physical chain
once array geometry and rating are declared \cite{P1}.

We adopt GHI as the primary scientific target, because it is the
quantity actually measured and its meaning does not depend on an
installation, and treat PV energy as a derived engineering output
obtained through an explicit, user-visible system model. Reporting
``kWh'' without a declared array size, tilt and period is not a
well-posed target, and we regard the separation as a correction rather
than a preference.

On models, the corpus converges on tree ensembles. Random forests
\cite{breiman2001} appear in \cite{P2,P4,P5}; boosted trees in
\cite{P2,P3,P4,P5}; and stacking \cite{wolpert1992} is the best overall
performer in \cite{P4}. Deep sequence models are explored in
\cite{P6}. We retain four estimators---random forest, histogram
gradient boosting \cite{ke2017}, extremely randomised trees
\cite{geurts2006}, and ridge regression \cite{hoerl1970} as a
deliberate linear floor---together with a stacked combination of the
four using a ridge meta-learner. Second representatives of a family
already present were withdrawn: a longer comparison table is not a
stronger result.

\subsection{Evaluation practice}

Our concern is with the evaluation layer, where the corpus is less
uniform.

\textbf{Splitting.} \cite{P2} uses a random split and describes it as
preventing leakage. For an autocorrelated series this places
observations minutes apart on both sides of the boundary. \cite{P5} is
explicit that its ``validation was performed using random sampling''
and names time-based splitting as future work. Section~\ref{sec:f2}
supplies that measurement.

\textbf{Target truncation.} \cite{P4} applies interquartile-range
outlier bounds of 165--331~W/m\textsuperscript{2} to solar radiation
and removes 10.16\,\% of observations. Clear-sky midday GHI at the
study's latitude exceeds 900~W/m\textsuperscript{2}, so the procedure
removes physically valid high-irradiance data and compresses the range
against which an RMSE of 47.30~W/m\textsuperscript{2} is subsequently
reported. We apply physical-range validation only.

\textbf{Metric definition.} The quantity defined as ``relative MAE'' in
\cite{P2} is $\mathrm{mean}(|O-F|/F)$, which is MAPE rather than MAE
normalised by the mean of observations---the convention that same work
uses for relative RMSE. We compute and label MAE, rMAE and MAPE
separately.

\textbf{Night hours.} Any model trivially predicts zero at night, and
retaining night hours inflates $R^{2}$ because between-group day/night
variance dominates the total sum of squares. \cite{P2,P3,P4} all
exclude night. We adopt the solar-zenith criterion of \cite{P3}
($\theta_{z} < 87^{\circ}$), the only physically defined rule among
them; fixed clock windows are latitude-dependent and fail at high
latitude.

These are observations about method, not allegations of error. Each is
a choice that a reader cannot presently price, which is what
Section~\ref{sec:results} sets out to change.

% =====================================================================
\section{Data and System}

\subsection{Data source}

All meteorological inputs are hourly ERA5 and ERA5-Land reanalysis
\cite{hersbach2020}, retrieved through the Open-Meteo Historical
Weather API \cite{openmeteo} without an API key. Retrieved variables
are GHI (\texttt{shortwave\_radiation}), direct normal and diffuse
irradiance, air temperature, relative humidity, dew point, surface
pressure, wind speed and direction, cloud cover and precipitation.

This is reanalysis, not pyranometer data: a modelled gridded product
representing a cell of order 9--25~km rather than a sensor at a
specific address. This bounds comparability---our figures are not
directly comparable to the ground-station measurements of \cite{P2} or
the on-site station of \cite{P4}---and we state it wherever results are
reported. Data are CC-BY~4.0; ERA5 is \copyright~ECMWF/Copernicus.

Experiments in Section~\ref{sec:results} use two to three years of
hourly data at Hyderabad, India ($17.385^{\circ}$N,
$78.487^{\circ}$E), typically ${\sim}26{,}000$ hours. Incomplete hours
are dropped and reported, never imputed; we do not manufacture
observations on which we then report metrics.

\subsection{Physical chain}

PV energy is obtained from irradiance through published models applied
in sequence: Erbs decomposition \cite{erbs1982} to separate diffuse and
direct components; HDKR transposition to the plane of array
\cite{duffie2013}; the Faiman cell-temperature model
\cite{faiman2008}; the PVWatts~v5 DC model \cite{dobos2014} with a
nameplate temperature coefficient; a multiplicative loss stack; and
inverter efficiency with AC clipping. Clear-sky reference irradiance
uses the Haurwitz model \cite{haurwitz1945}.

\subsection{Feature construction and the clear-sky index}

Rather than predicting GHI directly, models predict the clear-sky index
$k_{t} = \mathrm{GHI} / \mathrm{GHI}_{\mathrm{cs}}$, following
\cite{P1}. Dividing out deterministic geometry leaves the atmospheric
component---which is what a learned model can actually contribute---and
has a direct consequence for feature attribution that
Section~\ref{sec:f4} makes explicit.

Standardisation is fitted \emph{inside} each cross-validation fold on
training data only. Fitting a scaler on the full dataset before
splitting leaks test-set statistics, and does so in a way that no
subsequent metric will reveal.

% =====================================================================
\section{Methodology}

\subsection{Validation protocol}

Unless a comparison is explicitly about the split, all reported metrics
use a chronological split with an embargo between train and test
segments, and rolling-origin cross-validation. Night hours are excluded
by the $\theta_{z} < 87^{\circ}$ criterion in training and in every
reported metric.

\subsection{Metrics}

We report MAE, RMSE, rMAE and rRMSE (each normalised by the mean of
observations), MAPE with a zero-denominator guard, $R^{2}$,
Nash--Sutcliffe efficiency, and mean bias error. Probabilistic
forecasts additionally report pinball loss, empirical coverage (PICP),
normalised interval width (PINAW), and CRPS computed empirically from
quantiles \cite{gneiting2007}. Skill scores are computed against
persistence, clear-sky persistence, and hour-of-year climatology.

Train-set metrics are never reported alone, only beside validation
metrics, and a large gap is flagged. This is a direct response to
\cite{P2}, where distance-weighted $k$-NN attains a training relative
RMSE of 0.41\,\% against 5.77\,\% on test---train metrics for such a
model describe interpolation, not generalisation.

% =====================================================================
\section{Results}
\label{sec:results}

Each subsection isolates one methodological choice. Model, seed,
feature set and data are identical within each comparison; only the
disputed element varies.

% ---------------------------------------------------------------
\subsection{The interval-label convention}
\label{sec:f1}

Irradiance in an hourly archive is a mean over an interval, whereas
solar position is instantaneous. If position is evaluated at the
interval \emph{label} rather than at the interval's representative
instant, a systematic phase error follows. The convention is not
documented in a form we could rely on, so we measured it, sweeping
candidate offsets over a full year (Table~\ref{tab:f1}).

\begin{table}[!t]
\caption{Interval-offset sweep, Hyderabad, one year. The offset is
applied to the timestamp at which solar position is evaluated.}
\label{tab:f1}
\centering
\begin{tabular}{lrrr}
\toprule
Offset & Clear-sky & Apparent night-time & $\mathrm{corr}$(GHI, \\
(min) & exceedances & irradiance & clear-sky) \\
\midrule
$-60$ & 360 & 289 & 0.9340 \\
$\mathbf{-30}$ & \textbf{0} & \textbf{216} & \textbf{0.9490} \\
$0$   & 264 & 265 & 0.9372 \\
$+30$ & 697 & 428 & 0.8998 \\
\bottomrule
\end{tabular}
\end{table}

The $-30$~minute offset is the unique candidate producing zero
physically impossible clear-sky exceedances, and it simultaneously
maximises correlation with the clear-sky reference. We conclude that
hourly radiation is the mean over the \emph{preceding} hour and the
representative instant is 30~minutes before the label. Applying the
correction raised the platform's data-quality score from 97.7\,\% to
99.9\,\%.

Two points deserve emphasis. First, the naive choice---offset
zero---is not merely suboptimal but produces 264 physically impossible
observations per year, which any downstream physical-range check will
then either flag or silently absorb. Second, the sweep exposed a
genuine defect in our own implementation: the PV chain evaluated solar
position at the raw label while the feature pipeline used the
midpoint, a 30-minute inconsistency that displaced the day/night
boundary by a full hour at the terminator. It was caught by a test
asserting zero PV output at night. We report it because the class of
error is easy to introduce and invisible in aggregate metrics.

% ---------------------------------------------------------------
\subsection{Split optimism}
\label{sec:f2}

Table~\ref{tab:f2} varies only the splitting strategy.

\begin{table}[!t]
\caption{Effect of splitting strategy. Identical model, seed and
dataset; two years, Hyderabad.}
\label{tab:f2}
\centering
\begin{tabular}{lrr}
\toprule
Strategy & RMSE (W/m\textsuperscript{2}) & $R^{2}$ \\
\midrule
Chronological          & 74.68 & 0.9160 \\
Blocked random (days)  & 68.97 & 0.9376 \\
Random (hour level)    & 66.03 & 0.9427 \\
\bottomrule
\end{tabular}
\end{table}

Hour-level random splitting understates RMSE by 11.6\,\% and overstates
$R^{2}$ by $+0.027$. Blocking by whole days recovers roughly half the
gap, which is consistent with the mechanism: the leak is driven by
near-neighbour hours straddling the boundary, and blocking removes the
within-day cases while leaving day-to-day persistence intact.

An 11.6\,\% RMSE difference is comparable to the margin by which
competing models are separated in several published comparisons. Where
a study reports a random split, that margin is therefore not safely
interpretable as a model difference. This substantiates quantitatively
the concern that \cite{P5} raises about its own protocol.

% ---------------------------------------------------------------
\subsection{Prediction-interval calibration}
\label{sec:f3}

An interval that does not cover what it claims is worse than no
interval, because it will be relied upon. Quantile gradient boosting
trained directly for the 10th and 90th percentiles produced a nominal
80\,\% interval with \textbf{66.3\,\%} empirical coverage.

Applying conformalized quantile regression \cite{romano2019} restores
calibration (Table~\ref{tab:f3}).

\begin{table}[!t]
\caption{Empirical coverage of a nominal 80\,\% prediction interval,
before and after conformal calibration.}
\label{tab:f3}
\centering
\begin{tabular}{lr}
\toprule
Condition & Coverage (\%) \\
\midrule
Uncalibrated quantile GBM, all hours & 66.3 \\
\midrule
Conformalized, all hours   & 80.2 \\
\quad clear regime         & 80.3 \\
\quad cloudy regime        & 80.0 \\
\quad precipitation regime & 80.9 \\
\bottomrule
\end{tabular}
\end{table}

Coverage holds within two points inside every weather regime, using the
regime thresholds of \cite{P3}. Conformal calibration is not drawn from
the reconciled corpus and we label it an enhancement rather than an
implementation of prior work.

We also implemented and tested a seasonally stratified calibration
split. It did not robustly improve coverage: results were non-monotone
in its parameters (0.700, 0.725, 0.706, 0.790 as the embargo widened;
0.767 against 0.834 for six versus eight blocks). Non-monotone response
to a parameter that should act smoothly is instability, not signal, and
selecting the best-scoring configuration would have been tuning on the
test set. The simpler contiguous split was retained. We report the
negative result because the alternative---reporting only the
configuration that scored best---is precisely the practice that
produces irreproducible intervals.

% ---------------------------------------------------------------
\subsection{Feature importance depends on the target}
\label{sec:f4}

Grouped permutation importance, predicting the clear-sky index at
Hyderabad, is given in Table~\ref{tab:f4}.

\begin{table}[!t]
\caption{Grouped permutation importance when predicting the clear-sky
index. Shares of measured effect.}
\label{tab:f4}
\centering
\begin{tabular}{lr}
\toprule
Feature group & Share (\%) \\
\midrule
Temperature and moisture & 52.0 \\
Cloud and precipitation  & 21.0 \\
Solar geometry           & 17.0 \\
Time of year / day       & 6.0 \\
Pressure                 & 3.0 \\
Wind                     & 1.5 \\
\bottomrule
\end{tabular}
\end{table}

Solar geometry ranks third, whereas \cite{P5} finds geometry---angle of
incidence in particular---dominant. The results are consistent rather
than contradictory. \cite{P5} predicts raw PV power, in which geometry
is the largest single driver; we divide geometry out before fitting, so
what remains for the model to explain is the atmospheric component. In
that residual space humidity and cloud dominate, which agrees with
\cite{P2} finding temperature and humidity the leading meteorological
predictors.

The methodological point generalises: a feature-importance ranking is a
statement about a target, not about solar energy. Two studies reporting
different rankings may both be correct, and comparing their rankings
without first reconciling their targets is not meaningful.

% ---------------------------------------------------------------
\subsection{Baselines coincide at day-ahead horizons}
\label{sec:f5}

Naive persistence and clear-sky (``smart'') persistence score
143.4 and 143.3~W/m\textsuperscript{2} respectively at a 24-hour
horizon. This is expected rather than anomalous: solar declination
changes by at most ${\sim}0.4^{\circ}$ per day, so
$k_{t}(t-24\mathrm{h}) \cdot \mathrm{GHI}_{\mathrm{cs}}(t) \approx
\mathrm{GHI}(t-24\mathrm{h})$. The two diverge sharply at intra-day
horizons. It is also why baselines must be formed in physical units:
in clear-sky-index space they are algebraically the same function and
the comparison would be vacuous. A reported day-ahead improvement over
``smart persistence'' should therefore be checked against naive
persistence, which at that horizon is the same baseline.

% ---------------------------------------------------------------
\subsection{Scenario responses can conceal opposing pathways}
\label{sec:f6}

A perturbation scenario of $+6^{\circ}$C with halved wind speed
returned a net energy change of $+0.29\,\%$---counter-intuitive, since
hotter modules are less efficient. Decomposition
(Table~\ref{tab:f6}) shows two large opposing pathways.

\begin{table}[!t]
\caption{Decomposition of a $+6^{\circ}$C, halved-wind scenario.}
\label{tab:f6}
\centering
\begin{tabular}{lr}
\toprule
Pathway & Effect (kWh) \\
\midrule
Statistical (learned: warm hours tend to be sunny hours) & $+114.3$ \\
Physical (causal: hotter modules are less efficient)     & $-105.7$ \\
\midrule
Net & $+8.6$ \\
\bottomrule
\end{tabular}
\end{table}

The learned correlation and the causal physical response are of
comparable magnitude and opposite sign. Reporting the net figure alone
would conceal this, and would invite the reader to attribute a
near-zero response to insensitivity rather than to cancellation. We
surface both pathways separately. Doubling wind speed instead shows the
causal pathway dominating ($+43.2$~kWh thermal against $-7.0$~kWh
statistical), the physically expected result.

This is a general hazard for any scenario analysis conducted with a
model trained on observational data: the model has learned
correlations that a counterfactual intervention does not preserve.

% =====================================================================
\section{Discussion}

The effects reported here are not exotic. Each arises from a choice
that is ordinary, defensible in isolation, and rarely varied within a
single study. Their magnitudes, however, are on the same scale as the
differences by which published models are separated: an 11.6\,\% RMSE
shift from splitting alone, a 14-point coverage shortfall from omitting
calibration, and a feature ranking that reorders entirely with the
choice of target.

We therefore suggest three reporting practices, each cheap once
implemented. First, state the splitting strategy and, where a random
split is used, report the chronological figure alongside it; the
comparison costs one additional fit. Second, report empirical coverage
whenever a prediction interval is shown, since a nominal level is an
intention and not a measurement. Third, state the prediction target
precisely enough that a feature-importance ranking can be interpreted,
which for PV energy means declaring array size, tilt and period.

None of this requires agreement with our modelling choices. The
comparisons are internal: each holds the model fixed and varies only
the practice under examination.

% =====================================================================
\section{Limitations and Threats to Validity}
\label{sec:limits}

We state these plainly because they bound what the paper claims.

\textbf{Modelled, not measured.} No metered generation from an
installed system was available to us. Every PV energy figure is a
physical chain applied to reanalysis weather and has \emph{not} been
validated against real production. We consequently make no claim about
absolute forecast accuracy against physical plant, and readers should
not extract one. The findings in Section~\ref{sec:results} are
controlled internal comparisons, in which the absence of plant ground
truth affects both arms equally and therefore does not threaten the
contrast---but this is an argument for the specific claims made, not a
general licence.

\textbf{Reanalysis, not a pyranometer.} ERA5 represents a grid cell,
not a point. Local haze, dust and coastal cloud may differ materially
from the cell average, and our figures are not directly comparable to
studies using ground stations.

\textbf{Single-site experiments.} The measurements in
Section~\ref{sec:results} were conducted at one location. The
mechanisms are general---autocorrelation, interval labelling, quantile
miscalibration---but the effect \emph{sizes} should be expected to vary
with climate, and particularly with the frequency of rapid cloud
transitions. Multi-site replication is the most immediate extension.

\textbf{Interval-label finding is source-specific.} The $-30$~minute
result characterises this archive's convention. The measurement
procedure transfers to other sources; the offset should not be assumed
to.

\textbf{Scope.} Deep sequence models \cite{P6} and full copula-based
dynamic feature selection \cite{P3} are not implemented, and are
labelled as future work rather than approximated.

% =====================================================================
\section{Conclusion}

We reconciled six solar-forecasting studies into one platform and used
it to price three evaluation choices that the literature generally
leaves unpriced. Random splitting of autocorrelated hourly data
understates RMSE by 11.6\,\% and overstates $R^{2}$ by 0.027. Hourly
reanalysis irradiance is labelled at the end of its averaging interval,
and ignoring this produces 264 physically impossible observations per
year at a single site. A nominal 80\,\% prediction interval from
uncalibrated quantile regression covered 66.3\,\% of observations,
which conformal calibration restored to 80.2\,\% across every weather
regime. We further showed that permutation importance is a statement
about a target rather than about solar energy, reconciling two
published rankings, and that a scenario response can hide two large
opposing pathways.

Future work is field validation against metered generation, which is
the one thing that would let accuracy claims be made at all;
multi-site replication to establish how the effect sizes vary with
climate; and extension to the deep sequence models the corpus
identifies but this work does not implement.

% =====================================================================
\section*{Reproducibility}

The platform, the reproducibility manifest, and the interval
calibration script used for Table~\ref{tab:f1} accompany this
submission. All meteorological inputs come from a keyless public API,
so every experiment can be repeated at arbitrary coordinates without
credentials or a data-use agreement.

% =====================================================================
\section*{Acknowledgment}

The authors thank the Department of Computer Science and Engineering
for supporting this work. Weather and solar-resource data are provided
by Open-Meteo under CC-BY~4.0; ERA5 is \copyright~ECMWF/Copernicus.

% =====================================================================
\balance
\begin{thebibliography}{00}

\bibitem{P1} R.~Hobbs and S.~Joshi, ``Using open-source forecasts for
solar plant maintenance outage scheduling can reduce lost energy,''
\emph{IEEE J. Photovoltaics}, 2026, accepted for publication.

\bibitem{P2} T.~Mabodi and J.~Hammujuddy, ``Solar irradiance
forecasting for informed solar systems design and financing
decisions,'' \emph{SAIEE Africa Research Journal}, vol.~115, no.~3,
2024.

\bibitem{P3} Y.~Lyu and S.~Eftekharnejad, ``Probabilistic solar
generation forecasting for rapidly changing weather conditions,''
\emph{IEEE Access}, vol.~12, 2024.

\bibitem{P4} J.~Rosales Huamani \emph{et al.}, ``Efficient ML models
for solar radiation prediction using ensemble techniques: A case study
in low-rainfall arid climates,'' \emph{IEEE Access}, vol.~13, 2025.

\bibitem{P5} K.~Vijay Babu \emph{et al.}, ``Solar energy forecasting
using machine learning techniques for enhanced grid stability,''
\emph{IEEE Access}, vol.~13, 2025.

\bibitem{P6} M.~Hayajneh \emph{et al.}, ``Intelligent solar forecasts:
Modern ML models and TinyML role for improved solar energy yield
predictions,'' \emph{IEEE Access}, vol.~12, 2024.

\bibitem{hersbach2020} H.~Hersbach \emph{et al.}, ``The ERA5 global
reanalysis,'' \emph{Q. J. R. Meteorol. Soc.}, vol.~146, no.~730,
pp.~1999--2049, 2020.

\bibitem{openmeteo} P.~Zippenfenig, ``Open-Meteo.com weather API,''
2023. [Online]. Available: \url{https://open-meteo.com}

\bibitem{breiman2001} L.~Breiman, ``Random forests,'' \emph{Machine
Learning}, vol.~45, no.~1, pp.~5--32, 2001.

\bibitem{geurts2006} P.~Geurts, D.~Ernst, and L.~Wehenkel,
``Extremely randomized trees,'' \emph{Machine Learning}, vol.~63,
no.~1, pp.~3--42, 2006.

\bibitem{ke2017} G.~Ke \emph{et al.}, ``LightGBM: A highly efficient
gradient boosting decision tree,'' in \emph{Advances in Neural
Information Processing Systems}, vol.~30, 2017.

\bibitem{wolpert1992} D.~H. Wolpert, ``Stacked generalization,''
\emph{Neural Networks}, vol.~5, no.~2, pp.~241--259, 1992.

\bibitem{hoerl1970} A.~E. Hoerl and R.~W. Kennard, ``Ridge regression:
Biased estimation for nonorthogonal problems,'' \emph{Technometrics},
vol.~12, no.~1, pp.~55--67, 1970.

\bibitem{romano2019} Y.~Romano, E.~Patterson, and E.~Cand\`{e}s,
``Conformalized quantile regression,'' in \emph{Advances in Neural
Information Processing Systems}, vol.~32, 2019.

\bibitem{gneiting2007} T.~Gneiting and A.~E. Raftery, ``Strictly
proper scoring rules, prediction, and estimation,'' \emph{J. Amer.
Statist. Assoc.}, vol.~102, no.~477, pp.~359--378, 2007.

\bibitem{erbs1982} D.~G. Erbs, S.~A. Klein, and J.~A. Duffie,
``Estimation of the diffuse radiation fraction for hourly, daily and
monthly-average global radiation,'' \emph{Solar Energy}, vol.~28,
no.~4, pp.~293--302, 1982.

\bibitem{duffie2013} J.~A. Duffie and W.~A. Beckman, \emph{Solar
Engineering of Thermal Processes}, 4th~ed. Hoboken, NJ, USA: Wiley,
2013.

\bibitem{faiman2008} D.~Faiman, ``Assessing the outdoor operating
temperature of photovoltaic modules,'' \emph{Progress in
Photovoltaics}, vol.~16, no.~4, pp.~307--315, 2008.

\bibitem{dobos2014} A.~P. Dobos, ``PVWatts version~5 manual,''
National Renewable Energy Laboratory, Golden, CO, USA, Tech. Rep.
NREL/TP-6A20-62641, 2014.

\bibitem{haurwitz1945} B.~Haurwitz, ``Insolation in relation to
cloudiness and cloud density,'' \emph{J. Meteorology}, vol.~2, no.~1,
pp.~1--8, 1945.

\bibitem{pedregosa2011} F.~Pedregosa \emph{et al.}, ``Scikit-learn:
Machine learning in Python,'' \emph{J. Machine Learning Research},
vol.~12, pp.~2825--2830, 2011.

\end{thebibliography}

\end{document}
